% ============================================================
%  ABSTRACT
% ============================================================
Extratropical cyclones (ETCs) are recurring systems along the Southwestern Atlantic coast that generate winds and precipitation, whose spatial coincidence is usually assumed.
While the Northern Hemisphere has systematic studies on how these fields are distributed, the use of spatial density estimation and of univariate and multivariate Empirical Orthogonal Functions (EOF and MEOF) across the different phases of ETC development is more limited in the Southern Hemisphere.
This thesis develops such an integration to assess whether ETC dynamic intensity predicts the magnitude and location of its 10-meter wind and precipitation maxima, or whether both fields develop a distinct structural organization depending on the evolutionary phase.

The analysis is built on a cyclone-tracking catalog based on 850-hPa relative vorticity \citep{gramcianinov2020analysis} and the objective classification of the life cycle into four phases using \textit{CycloPhaser} \citep{coutodesouza2024new, deSouza2025CycloPhaser}, applied to ERA5 (1979--2020).
The sample was stratified by three cyclogenesis regions (SBR, LPB, ARG) and two intensity strata: the Global Population and the subset of the most intense cyclones (p90, drawn from the Global Population).
On this basis, probability density functions (PDF) of magnitude, spatial density estimation (PDFe) of location, univariate EOF decomposition with significance assessed via Bootstrap-$t$, and MEOF decomposition of the joint covariance were applied progressively, ultimately integrated into structural prototypes.

In the Global Population, a single EOF mode dominates wind variability (31--57\% of variance) and is correlated with vorticity; precipitation, in contrast, is spread across several modes (EOF1 only 11--16\%).
This organization changes in the p90 subset. Wind concentrates in the northwest quadrant (mode of 23--25~m\,s$^{-1}$, densities up to $27\times10^{-4}$ in SBR), while, from maturity onward, precipitation maxima move away from the center and form an arc-shaped structure to the east and southeast.
In the Global Population, MEOF1 captures 14.8--26.6\% of the joint covariance, with precipitation and wind gradients along perpendicular axes. In p90 it drops to 11.65--23.47\% and, at maturity, places both fields in opposite sectors of the cyclone, a separation supported by the significant correlation between the univariate PC1 of both fields ($|r|$ between $0.31$ and $0.45$), although wind dominates this joint mode in most cases.
The structural prototypes quantify this separation at over $8^{\circ}$--$9^{\circ}$ during maturity and decay, with ARG reaching the largest magnitude (${\sim}10$--$11^{\circ}$).
The magnitude of this separation varies regionally without a single identified factor, although its direction, precipitation to the east-southeast and wind to the northwest, remains consistent across the three regions.

These results show that, in the most intense ETCs of the Southwestern Atlantic, wind and precipitation maxima lie in opposite quadrants at maturity, and that central vorticity, which remains associated with the organization of wind, no longer describes that of precipitation in SBR and LPB.
This thesis combines PDFe, univariate EOF, and MEOF to quantify this separation by phase and genesis region in the Southwestern Atlantic, with implications for coastal risk in Brazil, Uruguay, and Argentina.

\vspace{1cm}

\noindent \textbf{Keywords:} South America; extratropical cyclones; wind at 10 meters; precipitation; spatial structure; multivariate EOF.
